% Transcription — fonds Alexandre Grothendieck, shelfmark 105, batch 2 (pages 21-40).
% Established from the facsimile archives/batches/105/batch-02.pdf (a mirror of
% https://grothendieck.umontpellier.fr/105.pdf), per the skill
% transcribe-grothendieck. The batch file carries no cover sheet: PDF page N
% is page 20+N of the folder (the Montpellier footer prints « 21 » ... « 40 »),
% and the archivists' pencil numbers agree wherever legible (« 21 », « 24 »,
% « 26 », « 27 », « 29 », « 31 », « 33 », « 35 », « 37 », « 39 »).
%
% Pass: Opus 5.5 (claude-opus-5-5), 2026-09-23 — first pass, unchecked against the pages by a human.
%
% Ten of the twenty pages are transcribed: 21, 24, 26, 27, 29, 31, 33, 35,
% 37, 39. Absent: pages 22, 25, 30, 32, 34, 36, 38 and 40, sheets of computer
% listing (banner pages, job logs, tables) carrying nothing of his — page 36
% has three words in pencilled capitals beside the listing (« UTDE UTIL
% UTSD »), a job-room hand, not mathematics; pages 23 and 28, blank sheets.
% The gaps in the numbering are the record.
%
% He paginates the run at the head of each leaf: page 21 is his 9, 24 his
% 10, 27 his « 10 bis », 29 his 11, 31 his 12, 33 his 13, 35 his 14, 37 his
% 15, 39 his 16; his pages 1-8 are presumably in batch 1 (not checked). Page 26 carries no number
% of his: a scratch sheet crossed out by two long pencil diagonals, between
% his 10 and 10 bis. A marginal note on page 33 sends the reader to his
% « p. 16 », i.e. page 39, which does state the result it announces.
% The whole batch is in English, as he wrote it; notes are in French.
% The inventory title « [Champs (stacks) 2] » is the archivists'; the first
% \section heading is ours (bracketed), the second, « Factorization », is his
% (the margin of page 33).
%
% What the batch is. Closed model categories and factorization systems, in
% his notation: A (later M) a category, Fl(A) its arrows; for Φ ⊂ Fl(A),
% Φ^* the arrows with the LLP with respect to Φ, Φ_* those with the RLP;
% (C, F, W) cofibrations, fibrations, weak equivalences, TC = C ∩ W,
% TF = F ∩ W.
%   21     (his 9) « Alternative presentation of theory »: a closed model
%          structure on a category with C = mono and W given; what remains
%          to prove: a) F^* = TC (a saturation property of W), b) TF = W ∩ F,
%          c) factorization (admitted), d) W stable under pull-back along F;
%          a retract argument for a).
%   24     (his 10) Left/right saturated classes, order-reversing bijection
%          Φ ↦ Φ_*, Ψ ↦ Ψ^*; factorization property of an orthogonal pair;
%          Prop.: under factorization, Φ = Ψ^* iff Φ stable by direct
%          factors; « Quillen pair », « Quillen triple (C, F, W) »; (M0),
%          (M1), (M2), (M5); « W strongly saturated ».
%   26     Scratch, cancelled: Kan complexes, Hom(I,Y) -> Y×Y; conditions
%          a) b) c) on W and the question whether they imply saturation
%          (fg, gf ∈ W ⇒ f, g ∈ W), with retract diagrams.
%   27     (his 10 bis) Equivalent conditions on a closed model category in
%          Quillen's sense; W = { pi | p ∈ C_*, i ∈ F^* }; cofibrant-fibrant
%          replacement diagram; W strongly saturated tested on M_cf.
%   29     (his 11) Constructing a Quillen triple from (C, W); Proposition:
%          (C, W) comes from a Quillen triple iff 1) C, W stable by direct
%          factors, W mildly saturated, 2) factorization for (C, C_*) and
%          (C ∩ W, (C ∩ W)_*), 3) W ⊃ C_*.
%   31     (his 12) 4) stability of W under base change by F and cobase
%          change by C (« strict Quillen triple »); NB on W as given by a
%          homological criterion versus the Kan-Quillen route via the
%          generating set (TC)_0; Gabriel-Zisman's anodyne extensions.
%   33-37  (his 13-15) Factorization: the small-object argument — amalgamated
%          sums X_D = X ⊔_A B over all squares (D), transfinite induction
%          X_α(f), Σ_α(f), accessibility cardinal π, the functorial
%          factorization id -> Σ; Example with Φ = saturation of Φ_0, and the
%          Zorn argument for d).
%   39     (his 16) Lemma: a left Q-saturated class contains isos, is stable
%          under composition, cobase change, amalgamated sums, ordinal direct
%          systems; Theorem ① ② for M stable by small direct limits and
%          accessible; « Next step is to try to get rid of » the smallness of
%          Φ_0.
%
% Notation fixed for the batch: his capital Φ is drawn most often in a
% cursive form close to a script L (𝔏) — page 24 alternates it with a
% round Φ in the same formulas — and is set \Phi throughout; Ψ as \Psi;
% the category of arrows Fl(M) as \mathrm{Fl}(M), the underlining he
% sometimes gives it (page 37) as \underline{\mathrm{Fl}}; the empty set he
% draws as φ (page 27) as \varnothing. Words he underlines are set \emph.
% The printed listing on page 27 carries in its margin the timestamps
% « 06/09/82-09:14:22 » and « CFT 1.09 (05/10/82) »; one note records it
% there, and nothing is inferred from it beyond what the paper shows.
\documentclass[11pt,a4paper]{article}
% Le préambule commun à toutes les transcriptions du fonds Grothendieck.
%
% Il tient en une idée : **l'incertitude se compose, elle ne se résout pas.**
% Une transcription qui lisse un mot douteux a détruit la seule chose pour
% laquelle on la faisait — un lecteur doit pouvoir savoir, à chaque endroit,
% ce qui était sur le feuillet et ce qui a été ajouté. Les six macros qui
% suivent sont donc l'appareil critique, et non de la décoration.
%
% Les mêmes commandes sont reconnues par `scripts/render.mjs`, qui produit la
% vue de lecture du site. Une macro ajoutée ici doit l'être aussi là-bas,
% faute de quoi le rendu échouera bruyamment — ce qui est voulu : un rendu
% silencieusement faux est pire qu'un rendu absent.

\ProvidesPackage{grothendieck}[2026/08/08 Transcriptions du fonds Grothendieck]

\usepackage{amsmath,amssymb,amsthm}
\usepackage{mathrsfs}
\usepackage{stmaryrd}
% Les diagrammes commutatifs sont partout dans ce fonds : c'est la langue
% même de la géométrie algébrique de Grothendieck, et une transcription qui
% les rendrait en prose ne transcrirait rien.
\usepackage{tikz-cd}
% Français par défaut — c'est la langue des feuillets — mais l'anglais est
% chargé aussi : la traduction et le résumé sont des documents anglais, et la
% typographie française (espaces avant les deux-points) y serait une faute.
% Ces documents basculent par \selectlanguage{english} en tête de corps.
\usepackage[english,french]{babel}
\usepackage{fontspec}
\usepackage{geometry}
\geometry{a4paper,margin=25mm,marginparwidth=28mm}
\usepackage{marginnote}
\usepackage{xcolor}
% ulem, not soul. `soul`'s \st and \ul are built on a character-by-character
% scan that XeLaTeX's font machinery breaks: the document fails with an
% undefined control sequence rather than degrading. ulem does the same two jobs
% and is Unicode-engine safe. `normalem` stops it hijacking \emph, which the
% transcriptions use for Grothendieck's own emphasis.
\usepackage[normalem]{ulem}
\usepackage{hyperref}
\hypersetup{colorlinks=true,linkcolor=black,urlcolor=black,pdfborder={0 0 0}}

% --- Métadonnées du lot -----------------------------------------------------
% Reprises telles quelles par le script de rendu, qui les affiche en tête de
% la vue de lecture. Elles fixent ce que le fichier prétend couvrir : sans
% elles, une transcription flotte sans point d'attache dans un fonds de seize
% mille feuillets.

\newcommand{\folder}[1]{\gdef\grothFolder{#1}}
\newcommand{\batch}[1]{\gdef\grothBatch{#1}}
\newcommand{\leaves}[2]{\gdef\grothFirst{#1}\gdef\grothLast{#2}}
\newcommand{\foldertitle}[1]{\gdef\grothFoldertitle{#1}}
\gdef\grothFolder{?}\gdef\grothBatch{?}\gdef\grothFirst{?}\gdef\grothLast{?}\gdef\grothFoldertitle{}

% --- Le résumé --------------------------------------------------------------
%
% Il ouvre la lecture modernisée, sous le titre « Résumé », et il est composé
% en petit. Ce n'est pas un ornement : c'est le seul passage du document qui
% ne soit pas des mathématiques, et un lecteur doit pouvoir le reconnaître
% avant de l'avoir lu — puis le sauter s'il connaît déjà le sujet, ou s'y
% arrêter s'il ne le connaît pas. Le filet qui le suit marque la reprise du
% corps.

\newenvironment{resume}
  {\small\setlength{\parskip}{0.45em}}
  {\par\medskip\hrule\medskip}

% --- Mots-clés ---------------------------------------------------------------
%
% En anglais, en clôture du résumé de la lecture modernisée : le vocabulaire
% moderne sous lequel chercher ce que les feuillets construisent. Le manifeste
% du site (`npm run manifest`) les extrait pour étiqueter la cote entière —
% cette ligne est la source unique des tags, il n'y a pas de fichier à côté.
\newcommand{\keywords}[1]{\par\smallskip\noindent
  {\footnotesize\textsc{Keywords} --- \textit{#1}}\par}

% --- L'appareil critique ----------------------------------------------------

% Le numéro de feuillet, en marge. C'est l'ancre qui permet de régler un
% désaccord sur une formule en regardant *un* feuillet plutôt qu'une chemise
% de six cents. Le site s'en sert aussi pour tourner les pages du fac-similé
% quand on fait défiler la transcription.
\newcommand{\leaf}[1]{%
  \marginnote{\footnotesize\color{gray!70}\textsf{#1}}%
  \label{leaf:#1}\ignorespaces}

% Illisible : on ne devine pas. Un mot inventé qui se lit comme les autres est
% le pire résultat possible.
\newcommand{\ill}{\textcolor{red!60!black}{[\,\ldots\,]}}

% Lecture douteuse : le mot est donné, et signalé comme incertain.
\newcommand{\uncertain}[1]{\uline{#1}}

% Ajout de l'éditeur — une lettre manquante, un mot sous-entendu.
\newcommand{\add}[1]{\textcolor{blue!50!black}{[#1]}}

% Biffé par Grothendieck lui-même. Ce qu'il a rayé fait partie du document :
% une rature dit souvent où la pensée a changé de direction.
\newcommand{\struck}[1]{\sout{#1}}

% Note du transcripteur, sur l'état matériel du feuillet.
\newcommand{\note}[1]{\par\smallskip{\small\color{gray!60!black}\textit{[#1]}}\par\smallskip}

% Note marginale de Grothendieck, distincte du corps du texte.
\newcommand{\marginal}[1]{\par\smallskip\noindent\hspace{1em}%
  {\small\color{gray!60!black}#1}\par\smallskip}

% --- Titre ------------------------------------------------------------------

\AtBeginDocument{%
  \begin{center}
    {\large\bfseries Fonds Alexandre Grothendieck --- cote n\textsuperscript{o}~\grothFolder}\\[2pt]
    {\itshape\grothFoldertitle}\\[4pt]
    {\small feuillets \grothFirst--\grothLast\ (lot \grothBatch)}\\[2pt]
    {\footnotesize\color{gray!60!black}Transcription établie d'après le fac-similé numérisé
     par l'Université de Montpellier.\\
     Les lectures incertaines sont soulignées, les illisibles notées [\,\ldots\,],
     les ajouts entre crochets.}
  \end{center}
  \vspace{1em}}

\endinput


\folder{105}
\batch{2}
\pages{21}{40}
\dating{[à partir de 1982]}
\watermark{Édition de démonstration}
\foldertitle{[Champs (stacks) 2] : notes manuscrites (s.d.).}

\begin{document}

\section{[Structures de modèles fermées, paires et triples de Quillen]}
\note{titre de l'éditeur, entre crochets : les feuillets n'en portent aucun.
Le texte est en anglais, tel qu'il l'écrit ; dans ce lot, la pagination de
sa main commence à 9.}

\page{21}
\note{p.\ 9 de l'auteur.}
Alternative presentation of theory, with possibly slightly less strong
result — as presentation of \uncertain{homotopy as conditions} on a $W$, so
that there should exist a closed model structure on $A$, s.th.\
$C = \text{mono}$, $W = $ the given ones, hence $TC = C \cap W$, $F$ defined
by \struck{lift.} RLP with respect to $TC = C \cap W$
\uncertain{(or by $(TC)_*$)}, $TF$ (« \uncertain{Kan-}\ill{} » maps) defined
by RLP with respect to $C$ ($TF = C_*$).

The problems which remain (as we know \uncertain{always}
$(TF)^* = (C_*)^* = C$)

ⓐ $F^* = \struck{A_{\text{any}}}\ ((TC)_*)^* = TC$ — and as \ill{}
$F^* \subset (TF)^* = C$, this amounts to: any monomorphism which has the LLP
with respect to $F$ is in $W$. (A saturation type property of $W$ with
respect to monomorphisms.)

ⓑ $TF = W \cap F$ — in terms of the assumptions already made previously upon
$W$, we know $TF \subset W$, i.e.\ $TF \subset W \cap F$, thus the problem is
$W \cap F \subset TF$, i.e.\ any $f \in W$ which has the RLP with respect to
$W \cap C$ is in $TF$ i.e.\ is \uncertain{Kan-}\ill{}.

c) Factorization $f = pi$ with either $p \in TF$, $i \in C$ (but this is
already known and independent of $W$), or $p \in F$, $i \in TC$
\add{and even $i \in C \cap W$} — and this shouldn't be too hard in terms of
3) and 4) by standard arguments. Thus I'll admit c for the time being.
\note{« and even $i \in C \cap W$ » est en interligne. Les conditions 3), 4)
et, plus bas, 5) sur $W$ ne sont pas posées dans ce lot.}

ⓓ Moreover, \struck{I'm} I'm interested in the following, quite important
for me, dual to \uncertain{(b)}: If $f \in F$, $X \to Y$, and
$(Y' \to Y) \in W$, then $(X' \to X) \in W$.

For handling ⓐ ⓑ ⓓ presumably I'll have to rely strongly upon factorization
c). Thus for \uncertain{a)}, using factorization of $f = pi$, and using
$f \in F^*$ we get a section $s$ of $p$ s.th.\ $sf = i$, thus $f$ is a
retract of $i \in W$, and \ill{} \ill{} if we assume 5).

\begin{tikzcd}
  & X' \arrow[d, "p \in F"] \\
  X \arrow[ur, "i \in W \cap C"] \arrow[r, "f"'] & Y \arrow[u, bend right=40, "s"']
\end{tikzcd}

\note{le petit diagramme est dans la marge inférieure, à gauche de la
dernière ligne. La lettre après « Thus for » est tracée comme le c) qui
précède ; mais l'argument (un monomorphisme ayant la LLP par rapport à $F$
est rétracte de $i \in W \cap C$) est celui qui répond à ⓐ.}

\page{24}
\note{p.\ 10 de l'auteur.}
$A$ category, if $\Phi \subset \mathrm{Fl}(A)$, let
\[
  \Phi^* = \{ u \in \mathrm{Fl}(A) \mid u \text{ has the LLP with respect to }
  \Phi \},
\]
\[
  \Phi_* = \{ u \in \mathrm{Fl}(A) \mid u \text{ has the RLP with respect to }
  \Phi \} .
\]
$\Phi$ is called left saturated if $\Phi = \Psi^*$ for some $\Psi$, i.e.\
$\Phi = (\Phi_*)^*$; $\Psi$ right saturated if \struck{it is}
$\Psi = \Phi_*$ for some $\Phi$ i.e.\ $\Psi = (\Psi^*)_*$. Left saturation,
right saturation, $\Phi \subset (\Phi_*)^*$, $\Psi \subset (\Psi^*)_*$,
associations $\Phi \mapsto \Phi_*$, $\Psi \mapsto \Psi^*$ are order
reversing bijections between left saturated and right saturated subsets.

\marginal{\uncertain{Right} saturation $\Rightarrow$ stability by
composition, \uncertain{base ch.}, \uncertain{retracts} ; dually for left
saturation.}
\note{note écrite en oblique dans la marge gauche, en regard des lignes
précédentes.}

The \emph{factorization property} for a pair $(\Phi, \Psi)$ of mutually
orthogonal subsets of $\mathrm{Fl}(A)$ ($\Phi \subset \Psi^*$ i.e.\
$\Psi \subset \Phi_*$), is the property \struck{that}
\[
  \forall f \in \mathrm{Fl}(A) \quad \exists \text{ factorization }
  f = pi, \text{ with } i \in \Phi,\ p \in \Psi .
\]
\struck{Let \ill{} be the left and right saturations respectively}
\struck{Prop.} \struck{Let}
\emph{Proposition} \add{(for orthogonal pairs $(\Phi, \Psi)$)}. Assume
factorization property holds. Then $\Phi = \Psi^*$ iff $\Phi$ stable by
direct factors. Dually $\Psi = \Phi_*$ iff $\Psi$ stable by retractions.
Hence $(\Phi, \Psi)$ are dual pair of subsets of $\mathrm{Fl}(A)$
(localizers) iff \struck{and} \struck{(they are orthogonal)} both $\Phi$ and
$\Psi$ are stable by direct factors.

\emph{Quillen pair} $(\Phi, \Psi)$: dual pair $+$ factorization
$\Leftrightarrow$ orthogonal pair $+$ factorization $+$ stability of
$\Phi$, $\Psi$ (by direct factors).

\emph{Quillen triple} $(C, F, W)$ \struck{Now let} with

a) $W$ \struck{\ill{}} \uncertain{weakly} saturated (isos in $W$,
\struck{and} and the \ill{} \ill{} \add{and stable by direct factors}) (M5)

b) $(C, F \cap W \overset{\text{def}}{=} TF)$ and
$(TC \overset{\text{def}}{=} C \cap W, F)$ are Quillen pairs i.e.\ (M1)
(M2) and stability of $C$, $F$, \add{$TC$, $TF$} \struck{under} under direct
factors (it is enough for $C$, $F$, provided it holds for $W$).
\[
  \begin{array}{ccc}
    C & TF & \\
    \cup & \cap & W \\
    TC & F &
  \end{array}
\]
\add{it would not be necessary here to assume the direct factor condition on
$W$, if we assume b)}
\note{cette insertion est reliée par une flèche au a).}

(hence $f \in W \Longleftrightarrow f = pi$, $i \in TC$,
\uncertain{$p$}${} \in TF$)

! This implies $W$ is strongly saturated i.e.\ if
$\gamma_W : A \to W^{-1}A$, $f \in W \Longleftrightarrow \gamma_W(f)$ iso.
\marginal{true at any rate if in $A$ finite direct \uncertain{limits} and
\uncertain{finite inverse limits} exist (M0)}

\page{26}
\note{feuillet de brouillon sans numéro de l'auteur, barré de deux longues
diagonales au crayon ; on le transcrit tel qu'il se lit.}
$Y$ is Kan. \qquad $X \to Y$ fibration, Kan.

$\underline{\mathrm{Hom}}(I, Y) \to Y \times Y$ is a Kan fibration (if $Y$
is Kan); is it a Serre fibration?
\[
  a \to Y \times Y, \quad \alpha, \beta \in Y_a, \quad
  \underline{\mathrm{Hom}}_{\alpha,\beta}(I_a, Y_a), \quad
  \underline{\mathrm{Hom}}_{\alpha_b, \beta_b}(I_b, Y_b)
\]

a) $\mathrm{Iso}_M \subset W$

b) If \ill{} \ill{} among $u$, $v$, $uv$ \uncertain{two} are in $W$, so is
the third

c) $W$ stable by retraction \struck{(stronger than \ill{})}

$\Rightarrow$ \uncertain{mild} saturation. But does it imply saturation,
namely $fg \in W$, $gf \in W \Longrightarrow f, g \in W$??

$A \underset{g}{\overset{f}{\rightleftarrows}} B$ \qquad
$fg \in W$, $gf \in W \Longrightarrow f, g \in W$?

$fg \in W_A$, $gf = \mathrm{id}_A \Longrightarrow f, g \in W$

\begin{tikzcd}
  A \arrow[r, "f"] \arrow[d, bend left=30, "i"] & B \arrow[d, bend left=30, "j"] \\
  A' \arrow[r, "f'"'] \arrow[u, bend left=30, "p"] & B' \arrow[u, bend left=30, "q"]
\end{tikzcd}

$f'i = jf$, \quad $fp = qf'$, \quad $pi = \mathrm{id}_A$, \quad
$qj = \mathrm{id}_B$, \quad $f' \in W$ $\overset{?}{\Longrightarrow}$
$f \in W$.

(i) $gf = \mathrm{id}_A$, $fg \in W \Rightarrow f, g \in W$

(ii) $gf \in W$, $fg \in W \Rightarrow f, g \in W$

(iii) $W$ stable \struck{by factors} direct factors

(ii) $\Rightarrow$ (i) $\Leftarrow$ (iii)

\note{autour de ces lignes, des esquisses : les carrés
$A \to A$, $A \xrightarrow{f} B$, $B \xrightarrow{fg} B$ exhibant $f$ comme
rétracte de $fg$ ; un encadré « a) b) c$_1$) (c$_2$) (c$_3$) » ; en bas à
gauche, $A \rightleftarrows B$ avec « $W$ », « ${}^{h}W$ »,
$gf \sim_h \mathrm{id}_A$ donc $gf \in W$, $fg \sim_h \mathrm{id}_B$ donc
$fg \in W$, et un triangle $X \leftarrow I \times X$ ; en bas à droite, un
schéma $X \to X$, $X \wedge X$ qu'on ne restitue pas.}

\page{27}
\note{p.\ « 10 bis » de l'auteur. Le listing sur lequel il écrit porte
imprimé dans sa marge « 06/09/82-09:14:22 » et « CFT 1.09 (05/10/82) ».}
$X \xrightarrow{f} Y$, \quad $X, Y \in M_{c,f}$, \quad $f \in F$.

Then $f \in TF \Longleftrightarrow$ retract property $\Longleftrightarrow$
i.e.

\[
  \begin{array}{cc}
    C & F \cap W = TF \\
    \cup & \cap \\
    TC = C \cap W & F
  \end{array}
\]
ⓐ $\begin{cases} C = (TF)^* \\ F = (TC)_* \end{cases}$
\qquad ⓑ $W = \bigl\{ pi \bigm| p \in \underbrace{C_*}_{\widetilde{TF}},\;
i \in \underbrace{F^*}_{\widetilde{TC}} \bigr\}$

$\widetilde{TF} \subset F$, $\widetilde{TC} \subset C$ (because $F$, $C$
\add{\uncertain{trivially}} saturated)

\struck{$TF = \widetilde{TF}$, $TC = \widetilde{TC}$ $\}$ because
$TF = F \cap W$, $TC = C \cap W$ are stable under direct factors [provided
we know $W$ is stable]}

we only have to \struck{prove} \uncertain{express}, to get \ill{} of
\ill{}, that
\[
  \boxed{\widetilde{TF} \subset W,\ \widetilde{TC} \subset W}
  \quad \text{and this is equivalent with} \quad
  \boxed{\widetilde{TF} = TF,\ \widetilde{TC} = TC}
\]
i.e.\ $TF$, $TC$ saturated i.e.\ $(C, TF)$, $(TC, F)$ Quillen pairs.

\emph{Equivalent conditions on a model category} (in Quillen's sense)
\struck{Closed model}

(i) satisfies ⓐ and ⓑ, i.e.\ closed model category in Quillen sense

$\Updownarrow$ (ii) $(C, TF = F \cap W)$ and $(TC = C \cap W, F)$ are Quillen
pairs

(iii) $C$, $F$, $W$ are stable under direct factors.

But I've to check that (i) (ii) (which are equivalent) does imply that
\emph{$W$} is stable under direct factors. But Quillen proves it even
implies $W$ is strongly saturated (i.e.\ $f \in W$ iff $\gamma(f)$ iso)
— at least under the assumption that finite direct and inverse limits exist
in $M$.

\uncertain{for any model category}

\begin{tikzcd}[column sep=small, row sep=small, nodes={font=\scriptsize}]
  & & X_2 \arrow[r, "C \text{ or } F", "g"'] & Y_2 \arrow[dr, "TF"] & & \\
  \varnothing \arrow[r, "C"] & X_1 \arrow[ur, "TC"] \arrow[dr, "TF"'] & & & Y_1 \arrow[r, "F"] & e \\
  & & X \arrow[r, "f"'] & Y \arrow[ur, "TC"'] & &
\end{tikzcd}

\note{deux flèches en pointillé, marquées l'une « $\varepsilon$ »,
l'autre « $F$ », vont de $X_1$ et de $X_2$ vers $Y_1$ ; on ne les restitue
pas. L'objet de départ est tracé $\varphi$ : on le lit comme l'objet
initial.}

$\gamma(f)$ iso $\Longleftrightarrow \gamma(g)$ iso; \quad
$f \in W \Longleftrightarrow g \in W$.

$W$ strongly saturated $\Longleftrightarrow$ for $f : X \to Y$,
$X, Y \in M_{cf}$, $f \in C$, $\gamma(f)$ iso $\Rightarrow f \in W$
$\Longleftrightarrow$ for $f : X \to Y$, $X, Y \in M_{cf}$, $f \in F$,
$\gamma(f)$ iso $\Rightarrow f \in W$.

\page{29}
\note{p.\ 11 de l'auteur.}
To construct any Quillen triple, we may proceed as follows: start with
\add{a localizer,} $C \subset \mathrm{Fl}(A)$, ⓐ $C$ stable \struck{\ill{}}
by direct factors, hence a pair $(C, C_* = TF)$ — we must check ⓑ
\emph{factorization condition for the pair} $(C, C_*)$. \struck{Hence} Then
$(C, TF)$ is a Quillen pair. Let next ⓒ $\boxed{W \supset TF}$ any
localizer which is \add{ⓓ \uncertain{mildly}} \emph{saturated and stable
under direct factor}. We now get $TC = C \cap W$ \add{$\subset C$}, which is
\struck{\ill{}} stable under direct factor. Take $F = (TC)_*$
($\supset TF = C_*$), we must check ⓔ \emph{factorization property for the
pair} $(TC, F)$ — thus we get that $(TC, F)$ is a Quillen pair. Lastly, we
must still check that $TF = F \cap W$, and by c) we know already
$TF \subset F \cap W$, we have to check still that $F \cap W \subset TF$.
But using a factorization $f = pi$, $p \in TF$, $i \in C$, as
$f, p \in W$ we get $i \in W$ hence $i \in W \cap C = TC$, hence there
exists $s : X' \to X$ s.th.\ $si = \mathrm{id}_X$ (retraction of $X'$ upon
$X$) and $fs = p$, hence $f$ a retraction of $p$, hence is in $TF$. OK.

\begin{tikzcd}
  X \arrow[r, "="] \arrow[d, "i"'] & X \arrow[d, "f \in F"] \\
  X' \arrow[r, "p"'] \arrow[ur, dashed, "s"] & Y
\end{tikzcd}

\note{à côté, un second schéma : $X' \rightleftarrows X$ par $s$ et $i$,
au-dessus de $Y \rightleftarrows Y$, avec $p$ et $f$ verticales — $f$
rétracte de $p$. Le $i$ du premier carré est marqué « $\in TC$ ».}

Thus

\emph{Proposition} In order for a pair $(C, W)$ to correspond to a Quillen
triple $(C, F, W)$ \add{fibrations}, it is n.s.\ that the following
conditions be satisfied:
\begin{enumerate}
  \item[1)] $C$, $W$ \emph{stable by direct factors}, $W$ mildly saturated
  \item[2)] Factorization conditions for $(C, C_*)$ and for
  $(C \cap W, (C \cap W)_*)$.
  \item[3)] $W \supset TF \overset{\text{def}}{=} C_*$
\end{enumerate}

We have to develop mainly a criterion to get factorization. [In the case of
interest to us, we get 1) free, 3) almost free (\ill{} comes in the
« test »-assumption though), and 2) \uncertain{causing} a little technical
problem.] I am however \ill{} factorization conditions, interested in the
\ill{} \ill{} in

\page{31}
\note{p.\ 12 de l'auteur.}
following properties (characterizing a \emph{strict} Quillen triple)

4) \struck{A any cobase change of} \add{$W$ stable by base} change by any
$g \in F$, and by cobase change by any $g \in C$.

\emph{NB} From my point of view, $W$ is given beforehand in a rather
tangible way, by a homological criterion, which turns out very handy. In
Kan-Quillen's point of view, $W$ was a little more hidden maybe in terms of
(say) \uncertain{ss} structures, they get to it via \ill{} the pair
$(C, TF)$ and $F \supset TF$ the \uncertain{Kan} fibrations, which is a
somewhat ad-hoc construction in terms of a \ill{} \uncertain{a-priori}
generating set $(TC)_0$ — corresponding roughly to expressing the
\uncertain{Kan} lifting property. What remains to be understood, in the
context of general test category, is the relationship between this and the
direct description $TC = C \cap W$, i.e.\ why this is just the
\emph{saturation} of $(TC)_0$. Maybe I should have a look at
Gabriel-Zisman and their treatment of « anodyne extensions ».

\section{Factorization}
\note{titre de sa main, souligné, écrit en oblique dans la marge supérieure
gauche du feuillet 33.}

\page{33}
\note{p.\ 13 de l'auteur.}
$M$ a category, $\Phi_0 \subset \mathrm{Fl}(M)$, $\Psi = (\Phi_0)_*$,
$\Phi_0 \subset \Phi \subset \overline{\Phi}_0 = \Psi^*$, we want to deduce
conditions on $\Phi_0$, $\Phi$ insuring that the pair $(\Phi, \Psi)$
satisfies factorization conditions, i.e.\
$\forall f \in \mathrm{Fl}(M)$, $\exists$ factorization $f = pi$,
$i \in \Phi$, $p \in \Psi$.
\add{NB.\ we don't use $\Phi \subset \uncertain{\overline{\Phi}_0}$}
\note{un astérisque est tracé au-dessus de $\Phi_0$ dans
$\Psi = (\Phi_0)_*$ ; sa portée n'est pas claire. L'insertion « NB.\ we
don't use … » est au-dessus de la ligne.}

\marginal{\ill{} the Problem \ill{} which \ill{} ; \uncertain{give the
result} \ill{} p.\ 16$^*$ \ill{}}
\note{note oblique dans la marge gauche, sous le titre ; « p.\ 16 » est sa
pagination et renvoie au feuillet 39, qui énonce le théorème.}

The idea is to take \add{(for given $f$)} \struck{all} \add{commutative}
diagrams

\begin{tikzcd}
  A \arrow[r, "u"] \arrow[d, "g"'] & X \arrow[d, "f"] \\
  B \arrow[r, "v"'] & Y
\end{tikzcd}

(D), with $g \in \Phi_0$, giving rise to

\begin{tikzcd}
  A \arrow[r, "u"] \arrow[d, "\alpha"'] & X \arrow[d, "i_D"] \arrow[dd, bend left=50, "f"] \\
  B \arrow[r] & X_D = X \sqcup_A B \arrow[d, "p_D"] \\
  & Y
\end{tikzcd}

\note{le carré supérieur est marqué « coc. » (cocartésien).}

and to define
\[
  \struck{\Sigma_1(f)}\; X_1(f) = \struck{\varinjlim}\;
  \coprod_{\substack{(X) \\ \text{all } D}} X_D
  \qquad \text{(amalgamated sum under } X)
\]
so that we get factorization

\begin{tikzcd}
  X \arrow[r, "i_1(f)"] \arrow[d, "f"'] & X_1(f) \arrow[dl, "\Sigma_1(f)"] \\
  Y &
\end{tikzcd}

with (hopefully) $i_1(f) \in \Phi$.
\note{le but de $i_1(f)$ est d'abord écrit $\Sigma_1(f)$, surchargé.}

For this first step, we need

a) $\Phi_0$ is small (we'll see later how to relax this condition)

b) Amalgamated sums indexed by $I$ with $\mathrm{card}\, I \leqslant
\mathrm{card}\, \Phi_0$ \struck{\ill{}} \add{or rather,
$\mathrm{card}\, \Phi_0 \times c$, where $c$ is
$\sup \mathrm{card}(\mathrm{Hom}(A, X) \times \mathrm{Hom}(B, Y))$} exist in
$M$, \struck{and}

c) $\Phi$ stable by \add{compositions, and by} cobase extensions, and by
filtering direct limits indexed by a set of cardinality
$\leqslant \mathrm{card}\, \Phi_0$ \struck{\ill{}} \add{rather
$\leqslant c$} \add{for $\alpha : A \to B$ in $\Phi_0$}.
\note{insertions et renvois très serrés en bout de lignes b) et c) ; leur
place exacte est incertaine.}

\marginal{\uncertain{do} \uncertain{with} the \ill{} \ill{}}

We then are making an transfinite induction to define $X_\alpha(f)$ and

\begin{tikzcd}
  X \arrow[r, "i_\alpha(f)"] \arrow[d, "f"'] & X_\alpha(f) \arrow[dl, "\Sigma_\alpha(f)"] \\
  Y &
\end{tikzcd}

$i_\alpha(f) \in \Phi$, with $\alpha$ any ordinal $< \alpha_0$ ($\alpha_0$
\struck{first} \add{\uncertain{suitable}} ordinal \struck{(\ill{})} to be
fixed later) \struck{define later in terms of \ill{}}:
\[
  \begin{cases}
    \Sigma_{\alpha+1}(f) = \Sigma_1(\Sigma_\alpha(f)), \quad
    X_{\alpha+1}(f) = \text{source } \Sigma_{\alpha+1}(f), \\
    i_{\alpha+1}(f) = \text{composition }
    X \xrightarrow{i_\alpha(f)} X_\alpha(f)
    \xrightarrow{i_1(\Sigma_\alpha(f))} X_{\alpha+1}(f) \\[1ex]
    \Sigma_\alpha(f) = \varinjlim_{\alpha' < \alpha} \Sigma_{\alpha'}(f)
    \quad \text{if } \alpha \text{ is a limiting ordinal.}
  \end{cases}
\]
\note{en regard de la première ligne, un schéma :
$f : X \to Y$, $\Sigma_\alpha(f) : X_\alpha(f) \to Y$,
$\Sigma_{\alpha+1}(f) = \Sigma_1(\Sigma_\alpha(f)) : X_{\alpha+1}(f) \to Y$.}

\struck{NB} We assume that $\Phi$ is stable too by the filtering direct
limits \struck{of} \struck{cardinality $\leqslant$} which occur in the second
transfinite induction step.

Let $\pi$ be a cardinal such that \struck{all objects of $\Phi_0$}
\add{for any $g \in \Phi_0$, source and target of $g$} are
$\pi$-accessible — i.e.\ $\mathrm{Hom}(A, -) : M \to \text{Sets}$ commutes
to filtering direct limits which are \struck{great wr} « large with resp.\
to $\pi$ ». We may take $\pi = \sup_{g \in \Phi_0} \pi_g$, where

\page{35}
\note{p.\ 14 de l'auteur.}
$\pi_g$ is the smallest cardinal with source and target of $g$ are
$\pi$-accessible.

\emph{Example} If the objects occurring as source and targets of $g$'s in
$\Phi_0$ are « of finite presentation », i.e.\ $\mathrm{Hom}(A, -)$ commutes
with any filtering direct limit, we get $\pi = \uncertain{0}$. In this case
we stop the induction with $\alpha_0 = $ first infinite ordinal,
\struck{the first} \struck{\uncertain{we mean}}. In case the objects $A$,
$B$ are not of finite presentation (but they are accessible) — we take the
smallest ordinal such that $\mathrm{Hom}(A, ?)$ commutes \uncertain{with}
direct limits indexed by \struck{ordi} ordered set \struck{great}
\uncertain{large} w.r.t.\ $\pi$. \add{Thus,} let $\alpha_0$ be the first
ordinal such the interval
$I_{\alpha_0} = \{ \text{all ordinals } \alpha \text{ with } \alpha <
\alpha_0 \}$ be \struck{great} larger w.r.t.\ $\pi$. We take the
factorization

\begin{tikzcd}
  X \arrow[r, "i_{\alpha_0}(f) = i(f)"] \arrow[d] & X_{\alpha_0}(f) = \underline{X}(f) \arrow[dl, "\Sigma_{\alpha_0}(f) = p(f)"] \\
  Y &
\end{tikzcd}

\note{le « $\pi = 0$ » est tracé d'un rond épais ; on ne décide pas s'il
s'agit de $0$, de $\omega$ ou d'autre chose. Dans le diagramme, un
$\Sigma$ est surchargé en $\underline{X}$.}

and we get that $i(f) \in \Phi$, $p(f) \in \Psi$, provided we make the
following assumptions:

d) The \struck{objects} source and target objects of the arrows
$g \in \Phi_0$ are \emph{accessible} (ok if $M$ a topos etc.\ —) and, if
$\pi_0$ is the cardinality of the smallest ordinal \ill{} such that for any
such $A$, $\mathrm{Hom}(A, -)$ commutes to the \struck{\ill{}} direct limits
of type $I_{\alpha_0}$, \struck{we assume} and \struck{$\pi = \sup(\pi_0$}
$\pi \geqslant \sup(\pi_0, \mathrm{card}\, \Phi_0)$ \add{rather, $c$}, we
assume that for any filtering direct systems with ordered set $I$ of
cardinality $\leqslant \pi$, and transitions maps in $\Phi$, the direct
limit exists and \struck{\ill{}} $Z_i \to Z$ are equally in $\Phi$.

Then we get, not only factorization $f = pi$, with $i \in \Phi$,
$p \in \Psi$, but even factorization $\underline{X}(f)$
\emph{functorial} with respect to variable $f$.
\note{les deux dernières lignes sont marquées d'un double trait en marge.}

\page{37}
\note{p.\ 15 de l'auteur.}
\emph{\ill{}} Assume \struck{\ill{}}

a) $\Phi_0$ is small, source and target of any $g \in \Phi_0$ are
accessible

b) $M$ stable by filtering direct limits, \struck{and by amalgamated}
\struck{\ill{}}

ⓒ \struck{Any} \add{For any} \struck{amalgamated} diagram
$A \xrightarrow{u} X$, $\alpha : A \to B$, with $\alpha \in \Phi_0$, the
amalgamated sum exists

\begin{tikzcd}
  A \arrow[r, "u"] \arrow[d, "\alpha"'] & X \arrow[d, "\alpha'"] \\
  B \arrow[r, "u'"'] & X'
\end{tikzcd}

and $\alpha' \in \Phi$.
\qquad Also c') $\Phi$ stable under compositions

ⓓ $\Phi$ « stable by \add{filtering} direct limits ».
\marginal{It is enough to know it w.r.t.\ ordered sets $I$ where sup of any
two elements exists}

Then $(\Phi, \Psi)$ satisfy the factorization condition, and \struck{more}
even, there exists a functor $\Sigma$

\begin{tikzcd}
  \underline{\mathrm{Fl}}(M) \arrow[rr, "\Sigma"] \arrow[dr, "t"'] & & \underline{\mathrm{Fl}}(M) \arrow[dl, "t"] \\
  & M &
\end{tikzcd}

and a functorial morphism \struck{\ill{}} in $[\mathrm{Cat}]_{/M}$
\add{possibly « large » categories over $M$}
\[
  \struck{\ill{}}\ \mathrm{id}_{\underline{\mathrm{Fl}}(M)} \to \Sigma
\]
with $i(f) \in \Phi$, $p(f) \in \Psi$.

\begin{tikzcd}
  X \arrow[r, "i(f)"] \arrow[d] & \underline{X}(f) \arrow[dl, "p(f)"] \\
  Y &
\end{tikzcd}

\emph{Example} Assume that $M$ is stable by amalgamated sums. Take
$\Phi = \overline{\Phi}_0 = \Psi^*$. Then c) \struck{and d)} is trivially
satisfied. What about d)? \uncertain{OK} \add{\uncertain{for compositions}}.
To get \struck{extension of} \add{the dotted $\varphi$}, we make
\uncertain{Zorn} \struck{\ill{}} \uncertain{induction} on pairs
$(J \subset I,\ u_J : A_J \to X)$ which make the diagram

\begin{tikzcd}
  A_{i_0} \arrow[r] \arrow[d, "u_0"'] \arrow[rr, bend left=30] & A_J \arrow[r] \arrow[dl, "u_J"] & B \arrow[d] \\
  X \arrow[rr, "f"'] & & Y
\end{tikzcd}

commutative, with \uncertain{obvious} order relation. \struck{This is
clearly}

\note{à gauche, le schéma qui pose le problème :
$A_{i_0} \to A_i \to B = \varinjlim A_i$, $u_0 : A_{i_0} \to X$,
$v : B \to Y$, $f : X \to Y$, et la flèche cherchée $\varphi : B \to X$ en
pointillé, marquée « induction ». La paire de la récurrence est écrite
$(J \subset I, u_J : \varinjlim_J A_j \to X)$, avec une surcharge.}
\[
  \mathrm{Hom}(B, X) \to \mathrm{Hom}(g_i, f) \quad \text{surjective}, \qquad
  \mathrm{Hom}(g, f) = \varprojlim \mathrm{Hom}(g_i, f)
\]

\page{39}
\note{p.\ 16 de l'auteur.}
\emph{Lemma} Let $\Phi \subset \mathrm{Fl}(M)$ be left Q-saturated (i.e.\
$\Phi = \Psi^*$ for some $\Psi$). Then \struck{$\Phi$ is stable by
compositions,}

a) $\Phi$ contains iso

b) $\Phi$ stable \struck{by} under composition and cobase changes
(\uncertain{including} \ill{})

c) $\Phi$ stable \struck{by} \add{under} \struck{direct} arbitrary
amalgamated sums

d) $\Phi$ stable \struck{by direct limits, indexed by any well ordered set
i.e.} \add{under ordinal direct systems, i.e.} \struck{if we have any
directed ordered set} if we have any well ordered set $I$ and
\struck{system} direct system $(A_i)_{i \in I}$ with
$g_{ij} : A_i \to A_j$ in $\Phi$, and
$A_{i_0} = \varinjlim_{j < i_0} A_j$ if $i_0 \in I$ is a limit, then
$\forall i_0 \in I$, $A_{i_0} \xrightarrow{g_{i_0}} B = \varinjlim A_i$ is
in $\Phi$.
\marginal{\emph{NB} a), b), d) $\Rightarrow$ c)}

\struck{The only}

Thus we get,

\emph{Theorem} Let $M$ be a category, stable by small direct limits, and
(\uncertain{weakly}) accessible. Then

① Let $\Phi_0 \subset \mathrm{Fl}(M)$\ \struck{\ill{}}$^*$, and
$\Phi_0 \subset \Phi \subset \overline{\Phi}_0 \overset{\text{def}}{=}
\Psi^*$ (with $\Psi = \Phi_{0*}$). If $\Phi$ \struck{stable and} satisfies
a) b) c) d)$^*$, then the $(\Phi, \Psi)$ satisfies factorization condition.
(Hence, if $\Phi$ stable by direct factors, $\Phi = \overline{\Phi}_0$.)

② (In order that $\Phi$ \add{(some data $\Phi_0$, $\Phi$)}
\uncertain{be saturated}, $\Phi = \overline{\Phi}_0$, it is n.s.\ that
$\Phi$ satisfy conditions a) b) c) d) and e) stability by direct factors.
In this case, factorization holds for the Q-orthogonal pair
$(\Phi, \Psi)$.

\marginal{We have to assume $\Phi_0$ small (\ill{}). We may \ill{} \ill{}
(\uncertain{which results} from a) b) d)) \uncertain{if} $\Phi_0$ \ill{}.
\uncertain{Alternatively} a) b$'$) c) d), where b$'$) is stability under
cobase \uncertain{changes} (without composition).}
\note{note oblique dans la marge gauche, en regard de ① et ② ; les
astérisques renvoient à cette note.}

\emph{NB.} In 2), \struck{besides} there was the preliminary assumption of
existence of $\Phi_0 \subset \Phi$ s.th.\ $\Phi_0$ and $\Phi$ \ill{} the
same left Q-saturation, with $\Phi_0$ \emph{small}. Next step is to try to
get rid of this.

\end{document}
